\section{Background and Related Work}

\paragraph{GUI visual search and scanpath prediction.}
GUI search models aim to predict how users look for interface targets such as buttons, links, controls, or text regions. Recent UI eye-tracking datasets and models have made this problem concrete for real screenshots, including UEyes for saliency across UI types, VSGUI10K for goal-directed visual search, and SeekUI for scanpath prediction from a GUI screenshot and text cue~\cite{jiang2023ueyes,putkonen2025vsgui,guo2026seekui}. Their evaluation is primarily target-present: the requested element is assumed to exist, and model quality is measured by scanpath or saliency agreement with human gaze. This leaves a gap for realistic GUI use, where the user may be on the wrong screen or the expected control may simply not be exposed.

\paragraph{GUI grounding and forced-choice evaluation.}
Many GUI grounding and agent benchmarks evaluate whether a model can identify an element matching a command or textual target. Recent multimodal UI models and GUI agents, such as Ferret-UI, ScreenSpot-Pro, and UI-TARS-style agents, emphasize screenshot-based perception and element grounding~\cite{you2024ferretui,li2025screenspotpro,wang2025uitars}. These settings often behave like forced-choice tasks: the model must select an element even when the command is underspecified, semantically broad, or absent from the screen. Our work studies what happens when this forced-choice assumption is applied to visual search. The key failure is not simply low localization accuracy, but overconfident present predictions on absent trials.

\paragraph{Cognitive stopping and target absence.}
Classical visual search includes both target-present and target-absent decisions. Models of attention and search describe how visual features, object-level binding, and guided selection shape the order in which items are inspected~\cite{treisman1980feature,wolfe2021guided}. Target-absent trials are not degenerate localization failures: they require continued search, reduced confidence, and a criterion for stopping. We use this cognitive framing as a lightweight post-hoc layer: if the generated scanpath does not accumulate evidence for a target-like candidate, and visible UI text also does not support the query, the system should be allowed to answer absent.

\paragraph{Abstention, open-set recognition, and VLM verification.}
The target-absent setting is related to selective prediction and open-set recognition: a model should sometimes reject an input instead of returning an unsupported label~\cite{geifman2017selective,scheirer2013openset}. Standard abstention methods usually operate on classifier confidence or distributional novelty. In contrast, GUI search gives us behavioral and interface-specific evidence: where the predicted path goes, what candidate text is visible, and whether the screenshot contains a plausible target. OCR provides an interpretable but brittle signal for text targets; VLMs provide broader visual reasoning but can be prompt-sensitive. Our contribution is to connect these threads by treating target absence as a search-stopping and verification problem grounded in scanpath and UI evidence.
